\section{Conclusion}\label{sec:conclusion}
% Drafting notes (not typeset): the only outcome-dependent statement is the \acc{} slot in the first
% paragraph (grammar: header of 11_replication.tex). Nothing else here may depend on the acceptance run.

A model comparison is only as credible as its weakest link: the sampler, convergence at every
origin, the Monte Carlo error of the score, the variable ordering, the validity of the test, and
choices made after the scores were seen. \pkg{rai-bvar} makes each of these links checkable. The
econometrics is not new; the contribution is the integration, the contracts between the components,
and the evidence, with its measured power and its stated gaps, that the samplers target the intended
posteriors. The validation code behind that evidence ships with the library. % [V] src/rai_bvar/validation/ (part of the package; dossier 01 §2, commitment 4)
In the pre-registered replication of \citet{clark2015macroeconomic}, the acceptance run's outcome
by criterion is \acc{criteria[*].met}. \todo{One clause stating which criteria the
acceptance run met and, if any failed, which, from \code{criteria[*].met}; no further
outcome-dependent claim.}

\paragraph{A checklist.}
Lessons for forecasters and for implementers of the same algorithms:
\begin{enumerate}[leftmargin=*,itemsep=2pt,topsep=4pt]% as the list of Section 4
  \item Write every kernel as a pure function of injected noise, exposing each exact draw's
    conditional and each Metropolis--Hastings log acceptance ratio (\cref{sec:arch-mcmc}).
  \item Write the reference samplers, simulators and joint densities independently of the kernels,
    and enforce the separation with an import test (\cref{sec:validation-harness}).
    % [V] tests/validation/test_log_joint.py::test_references_never_import_kernels
  \item Check an identity at a joint draw for each conditional and Metropolis--Hastings step; run
    prior-started Geweke chains and SBC on data-dependent quantities such as the log-likelihood
    (\cref{sec:validation}).
  \item Measure the suite's power with injected bugs, among them errors of published samplers,
    and degrade an approximation on purpose to test its exact correction
    (\cref{sec:validation-mutation}).
    % [V] tests/mutation/mutants.py: M1 (CCM-2019 triangular step), M2 (Primiceri 2005 Alg. 1 indicator order); M7 with the COARSE1 table
  \item Treat numerical failures as values: never add jitter; fix the kernel, not the test prior
    (\cref{sec:robustness}).
  \item Gate verdicts on convergence at every origin, and size the MCMC budget with fit-only
    diagnostics before any score exists (\cref{sec:eval-verdict,sec:replication-prereg}).
  \item Do not read the delta-method MCSE of an RB log score as a bound on its tail error: inspect
    the effective number of weights \citep{kish1965survey}; at $h > 1$, integrate the path out and
    use enough forecast draws (\cref{sec:eval-tail}).
  \item Report RB and Gaussian-approximation scores side by side, the ordering spread next to the
    gain, and an order-invariant negative control (\cref{sec:replication-design,sec:replication-results}).
  \item Treat Diebold--Mariano $p$-values under expanding windows as descriptive, and state which
    route produced each verdict (\cref{sec:eval-tests}).
  \item Pre-register settings as code, say how blind the pre-registration is, and check the absence
    of look-ahead at run time (\cref{sec:replication-data,sec:replication-prereg}).
\end{enumerate}

\paragraph{Roadmap and availability.}
The roadmap continues by phase (\cref{sec:limitations}); we make no promise about results. % [R] SPEC.md §13
The next priority is the size and power study of the verdict engine on simulated data, followed by
the Phase~3 prior axis, so that a two-axis ablation can be shown. \pkg{rai-bvar} is licensed under
Apache-2.0.
% [V] LICENSE; pyproject.toml (license = "Apache-2.0")
\todo{TODO 20: repository URL, release tag and commit hash, Zenodo DOI for the release and the
replication outputs, and the single command that rebuilds every table and figure
(the \code{Makefile} in \code{docs/paper}).}
